%%
%% spring-lecture01.tex
%% 
%% Made by Alex Nelson <pqnelson@gmail.com>
%% Login   <alex@lisp>
%% 
%% Started on  2026-04-04T08:31:56-0700
%% Last update 2026-04-04T08:31:56-0700
%% 

\lecture{}

\begin{node}
The subject of this course is associative algebras.
The textbook for the class is Richard S.~Pierce's \textit{Associative Algebras}
(Springer, 1982)~\cite{pierce1982associative}.
\end{node}

\begin{node}
Unless otherwise stated, $R$ is a commutative unital ring.
\end{node}

\begin{definition}
An \define{$R$-Algebra} (also called an ``algebra over $R$'') is a
unital right $R$-module $A$ equipped with an associative bilinear
product
\begin{equation}
\begin{split}
&A\times A\to A\\
&(x,y)\mapsto xy
\end{split}
\end{equation}
specifically such that
\begin{enumerate}
\item Associativity: $\forall x,y,z\in A\ldotp(xy)z=x(yz)$
\item Left distributive over addition: $\forall x,y,z\in A\ldotp x(y+z)=xy+xz$
\item Right distributive over addition: $\forall x,y,z\in A\ldotp (x+y)z=xz+yz$
\item Compatibility with scalar multiplication: $\forall x,y\in A\forall r\in R\ldotp (xy)r=x(yr)=(xr)y$
\item Unitality: $\exists 1_{A}\in A\ldotp\forall x\in A\ldotp 1_{A}x=x1_{A}=x$
\end{enumerate}
Informally ``algebra = ring + module''.
\end{definition}

\begin{note}
The ring morphism $R\to Z(A)$ sending $a\mapsto 1_{A}a$ determines the
structure of an $R$-algebra on $A$.
\end{note}

\begin{note}
Since $R$ is commutative, we define $rx:=(1_{A}r)x=(1_{A}x)r=xr$,
which makes $A$ into a left $R$-module, so $(r_{1}r_{2})x=r_{1}(r_{2}x)$.
\end{note}

\begin{node}
If $A$ is an $R$-algebra and $M$ is a right (resp., left) $A$-module, 
then $M$ is a right (resp., left) $R$-module by $mr:=m(1_{A}r)$
(resp., $rm:=(1_{A}r)m$).
\end{node}

\begin{remark}
For an example of a nonassociative unital algebra, the Octonions.

For an example of a nonassociative nonunital algebra, any Lie algebra.
\end{remark}

\begin{definition}
Let $A$ and $B$ be $R$-algebras. A \define{$R$-Algebra Morphism}
from $A$ to $B$ is a $R$-linear function $\phi\colon A\to B$ such that
\begin{enumerate}
\item Multiplicative: $\forall x,y\in A\ldotp\phi(xy)=\phi(x)\phi(y)$.
\end{enumerate}
\end{definition}

\begin{definition}
For $R$-algebras $A$ and $B$, we define an \define{$A$-$B$-Bimodule}
$M$ is simultaneously a left $A$-module and a right $B$-module such
that
\begin{enumerate}
\item $\forall a\in A\ldotp\forall b\in B\ldotp\forall u\in M\ldotp (xu)y=x(uy)$
\item $\forall u\in M\ldotp\forall r\in R\ldotp rm=mr$ --- i.e., $(1_{A}r)m=m(1_{B}r)$
\end{enumerate}
We follow Pierce and call use the term \define{$A$-Bimodule} as a
synonym for an $A$-$A$-Bimodule.
\end{definition}

\begin{node}
For more conventions, see Pierce~\cite[\S1.1]{pierce1982associative}.
\end{node}

\begin{definition}\label{defn:1.2}
Let $G$ be a monoid. We denote by $R[G]$ the \define{Monoid Algebra}
of $G$ over $R$ consisting of the set
\begin{equation}
R[G]=\{\xi\in R^{G}=\Map(G,R)\mid\xi(x)=0\mbox{ for almost all }x\in G\}
\end{equation}
(where $\Map(G,R)$ is the collection of set-theoretic functions from the set
underlying $G$ to the set underlying $R$) equipped with
\begin{enumerate}
\item addition defined by: $\forall\xi,\eta\in R[G]\ldotp\forall a,b\in R\ldotp (\xi a+\eta b)(x)=\xi(x)a+\eta(x)b$
\item multiplication: $\forall\xi,\eta\in R[G]\ldotp (\xi\eta)(x)=\sum_{\substack{y,z\in G\\ yz=x}}\xi(y)\eta(z)$
\end{enumerate}
We usually just write an element $\xi\in R[G]$ as something like
\begin{equation}
\xi = \sum_{g\in G}g\xi_{g}
\end{equation}
where only finitely many of the coefficients $\xi_{g}\in R$ are
nonzero (so the sums all make sense).
\end{definition}

\begin{lemma}\label{lemma:1.3}
For a (possibly non-associative) $R$-algebra $A$, if $X\subset A$
generate $A$ as an $R$-module and the associative (resp., unital) law
holds in $X$, then $A$ is associative (resp., unital).
\end{lemma}

\begin{proof}
It follows directly from $A=\sum_{x\in X}xR$,
since
\begin{subequations}
  \begin{align}
\left((\sum_{i}x_{i}a_{i})(\sum_{j}x_{j}b_{j})\right)(\sum_{k}c_{k}x_{k})
&=\sum_{i,j,k}(x_{i}x_{j})x_{k}(a_{i}b_{j}c_{k})\\
&=\sum_{i,j,k}x_{i}(x_{j}x_{k})(a_{i}b_{j}c_{k})\\
&=(\sum_{i}x_{i}a_{i})\left((\sum_{j}x_{j}b_{j})(\sum_{k}c_{k}x_{k})\right).
  \end{align}
\end{subequations}
Hence the claim for associativity. The argument for unital is similar.
\end{proof}

\begin{proposition}\label{prop:1.4}
\begin{enumerate}
\item\label{item:1.4:1} $R[G]$ is an $R$-algebra
\item\label{item:1.4:2} As an $R$-module, $R[G]$ is free with basis $G$
\item\label{item:1.4:3} If $A$ is an $R$-algebra, and if $\phi\colon G\to A$ is a module
  morphism, then $\phi$ extends uniquely to the $R$-algebra morphism
  $\widetilde{\phi}\colon R[G]\to A$.
\end{enumerate}
\end{proposition}

\begin{proof}
\ref{item:1.4:1} and \ref{item:1.4:2}: Checking multiplication is a
closed bilinear map is routine. For $x\in G$, set $\chi_{x}\in R[G]$
such that
\begin{equation}
\chi_{x}(y)=\delta_{x,y}=\begin{cases}1 & x=y\\
0 & \mbox{otherwise}
\end{cases}
\end{equation}
then these form a basis for $R[G]$. Hence $R[G]$ is free and its basis
$\{\chi_{g}\mid g\in G\}=X$. (We need to check associative and unital
laws hold for $X$ to apply Lemma~\ref{lemma:1.3}.) Moreover, we have
\begin{equation}
\chi_{x}\chi_{y}=\chi_{xy}
\end{equation}
[Pf: $(\chi_{x}\chi_{y})(z)=\sum_{z_{1}z_{2}=z}\chi_{x}(z_{1})\chi_{y}(z_{2})=\delta_{z,xy}$]
and also $\chi_{1}$ is a unity in $X$. So by Lemma~\ref{lemma:1.3},
$R[G]$ is an $R$-algebra, where $\chi_{1}=1_{R[G]}$.

\ref{item:1.4:3} For any $R$-algebra $A$ and any $R$-algebra morphism
\begin{equation}
\widetilde{\phi}\colon R[G]\to A
\end{equation}
is uniquely determined by $\widetilde{\phi}|_{G}$. So we may take 
\begin{equation}
\widetilde{\phi}|_{G}=\phi.
\end{equation}
By $\phi(xy)=\phi(x)\phi(y)$ and $\phi(1)=1$, we see $\phi$ is a
monoid morphism. Hence $\widetilde{\phi}$ is a Ring morphism.
\end{proof}

\begin{example}\label{ex:1.5}
Let $X$ be a set of symbols (e.g., $X=\{x_{0},x_{1},\dots\}$, possibly
finite or infinite). Let $G_{X}$ be the free monoid on $X$. By
Proposition~\ref{prop:1.4}, every mapping $X\to A$ uniquely extends to
$R[G_{X}]\to A$. If we replace $G_{X}$ by the free \emph{commutative}
monoid $H_{X}$, then $R[X]:=R[H_{X}]$ is the ring of polynomials with
unknowns in $X$.
\end{example}

\begin{definition}\label{defn:1.6}
\begin{enumerate}
\item 
Let $A$ be an $R$-algebra. Let $M$ and $N$ be right (resp., left)
$A$-modules. We define the \define{Module of $A$-module morphisms from
  $M$ to $N$}, denoted $\hom_{A}(M,N)$, for the $R$-module of all
$A$-module morphisms where addition is given pointwise
\begin{subequations}
\begin{equation}
(\phi+\psi)(u)=\phi(u)+\psi(u)
\end{equation}
and (right) scalar multiplication is given by
\begin{equation}
(\phi a)(u)=\phi(u)a.
\end{equation}
\end{subequations}
\item 
When $M=N$, then $\hom_{A}(M,N)=:\End_{A}(M)$ is the
\define{Endomorphism Algebra} of $M$ with multiplication $(\phi\psi)(u):=(\phi\circ\psi)(u)=\phi(\psi(u))$
\end{enumerate}
\end{definition}

\begin{note}
\begin{enumerate}
\item For any right $A$-module $M$, we see $\End_{A}(M)$ has an $A$-bimodule
structure. 
\item For any left $A$-module $M$, we see $\End_{A}(M)$ has an
$A$-bimodule structure.
\item If $A$ and $B$ are $R$-algebras and if $M$ is an
  $A$-$B$-bimodule, then for any $x\in A$ we define
  \begin{equation}
\lambda_{x}(u):=xu
  \end{equation}
  and $\lambda_{x}\in\End_{B}(M)$ [Pf: for any $y\in B$, $u\in M$,
    $x\in A$, we have $\lambda_{x}(uy)=x(uy)=(xu)y=\lambda_{x}(u)y$].
  Similarly,
\begin{equation}
\rho_{y}(u):=uy
\end{equation}
is in $\rho_{y}\in\End_{A}(M)$ for any $y\in B$. Also
\begin{equation}
\begin{split}
\lambda\colon &A\to\End_{B}(M)\\
&x\mapsto\lambda_{x}
\end{split}
\end{equation}
[Pf: show $\lambda_{xr}=\lambda_{x}r$ since $\lambda_{xr}(u)=(xr)u=x(ru)=(xu)r=\lambda_{x}(u)r$].
Similarly,
\begin{equation}
\begin{split}
\rho\colon&B\to\End_{A}(M)\\
&y\mapsto\rho_{y}
\end{split}
\end{equation}
is an $R$-algebra \underline{anti}morphism (i.e., a morphism $\op{B}\to\End_{A}(M)$).
\end{enumerate}
\end{note}

\begin{node}
If $M$ is a left $A$-module, then $M$ is an $A$-$R$-bimodule (with
$R$-action given by $m\cdot r=(1_{A}r)m$). We have a bijection
\begin{equation}
\begin{split}
\{\mbox{left $A$-modules}\}&\stackrel{1:1}{\longleftrightarrow}\{\mbox{representation of $A$}\}\\
M&\mapsto\lambda\\
[M_{\phi}]&\mapsfrom\phi
\end{split}
\end{equation}
where $[M_{\phi}]$ is defined by $x\cdot u=\phi(x)u$, and a
\define{Representation} of an $R$-algebra $A$ is an $R$-module $M$
equipped with an $R$-algebra morphism $A\to\End_{A}(M)$.
\end{node}

\begin{definition}
Let $\phi,\psi\colon A\to\End_{A}(M)$ be representations of $A$ on
$M$. We say $\theta\in\End_{R}(M)$ is an \define{Intertwiner} of
$\phi$, $\psi$ if $\theta\phi(x)=\psi(x)\theta$.
\end{definition}